% Systems / Rust / open source version: Canonical, Red Hat, Collabora,
% JetBrains, Cloudflare, Igalia, System76, Ferrous Systems.
% Build: pdflatex cv-systems-rust.tex (twice). Needs the lato package.
\documentclass[10pt]{article}
\usepackage{cv-style}
\hypersetup{pdftitle={Abdelrahman Khalifa - Systems Software Engineer - CV},pdfauthor={Abdelrahman Khalifa}}

\begin{document}

\cvheader{Abdelrahman Khalifa}{Systems Software Engineer --- Rust $\cdot$ C $\cdot$ Linux $\cdot$ GSoC 2026 (Linux Foundation)}

\section*{Summary}
Systems software engineer working in Rust, C and modern C++ on Linux. Completed Google Summer of Code 2026 with The Linux Foundation (OpenPrinting), building the COSMIC desktop's printer stack and landing 7 PRs upstream in cups-rs and libcups. Comfortable with FFI, IPC, async code and maintainer code review.

\section*{Open Source}
\cventry{Google Summer of Code 2026 Contributor --- The Linux Foundation / OpenPrinting}{May 2026 -- Oct 2026}
\cvsub{COSMIC desktop printer setup tool in Rust; design discussions, code review and documentation with OpenPrinting and System76 mentors \enspace|\enspace \cvlink{https://github.com/Abd002/GSOC-2026}{github.com/Abd002/GSOC-2026} \enspace|\enspace \cvlink{https://github.com/Abd002/cosmic-printers}{github.com/Abd002/cosmic-printers}}
\begin{itemize}
  \item Designed a printer-management stack as 5 Rust crates: \texttt{printers-core} (types, grouping queues into physical devices), \texttt{-client} (async varlink client), \texttt{-server} (CUPS, IPP, DNS-SD and Printer Application backend), \texttt{-ui} (iced/libcosmic) and \texttt{-app}.
  \item Exposed the backend over varlink IPC from \texttt{cosmic-settings-daemon}, so the C FFI code never loads in the Settings GUI process; the standalone app falls back to running the server in-process.
  \item Built a Printers page and an Add Printer flow for COSMIC Settings: live updates, print queue control, supply levels, and legacy printers set up through Printer Applications (PAPPL) instead of PPD drivers.
  \item Tested on real HP printers and Printer Applications (Ghostscript, Gutenprint, HPLIP); CI checks formatting, Clippy and tests on several architectures and runs CUPS, Avahi and \texttt{ippeveprinter} to test discovery.
\end{itemize}
\cventry{Upstream contributions --- OpenPrinting cups-rs (Rust) and libcups (C)}{Jan 2026 -- Oct 2026}
\begin{itemize}
  \item cups-rs: ported to libcups3 with CUPS 2 and CUPS 3 builds and CI; added IPP management operations, Printer Application support, connections by IP/port, \texttt{lpoptions} support and FFI bindings for the libcups3 DNS-SD API (\cvlink{https://github.com/OpenPrinting/cups-rs/pull/18}{PRs \#18--\#22}, merged).
  \item libcups: fixed crashes in \texttt{cupsDNSSDDelete()} and a leaked client in the Avahi DNS-SD code (\cvlink{https://github.com/OpenPrinting/libcups/pull/164}{\#164}, \cvlink{https://github.com/OpenPrinting/libcups/pull/165}{\#165}, merged).
  \item COSMIC: Printers page in Settings (\cvlink{https://github.com/pop-os/cosmic-settings/pull/2182}{cosmic-settings \#2182}) and the varlink printers server (\cvlink{https://github.com/pop-os/cosmic-settings-daemon/pull/188}{cosmic-settings-daemon \#188}), in review.
  \item Winter of Code 2026: implemented the \texttt{lp} and \texttt{lpstat} commands in Rust on cups-rs (\cvlink{https://github.com/Gmin2/cups-rs/pull/11}{PR \#11}).
\end{itemize}
\cventry{Contributor --- Boa JavaScript engine (Rust), Rust Clippy, Zephyr RTOS}{Mar 2024 -- Dec 2025}
\begin{itemize}
  \item Boa: optimized \texttt{JsStr::code\_points()} for Latin1 strings, avoiding UTF-16 conversion; with tests (\cvlink{https://github.com/boa-dev/boa/pull/4553}{PR \#4553}).
  \item Rust Clippy: fixed \texttt{ref\_as\_ptr} false positives for non-temporary references in \texttt{let}/\texttt{const} (\cvlink{https://github.com/rust-lang/rust-clippy/pull/16214}{PR \#16214}).
  \item Zephyr RTOS: implemented the \texttt{IP\_MULTICAST\_LOOP} socket option across 6 networking modules, merged upstream (\cvlink{https://github.com/zephyrproject-rtos/zephyr/commit/b11703623c02f5afa4a5b0b7abdce3e324d7d470}{commit b117036}).
\end{itemize}

\section*{Experience}
\cventry{Embedded Systems Engineer --- EKSON}{Aug 2024 -- Present}
\begin{itemize}
  \item Built a modern C++/Qt control application for motion simulators and optimized its performance by \textbf{\textasciitilde40\%} through asynchronous processing and fewer threads; reliable UART control.
  \item Developed Zephyr RTOS firmware for an ESP32-S3 3-DoF simulator, migrating from FreeRTOS for more deterministic behavior and production readiness.
\end{itemize}
\cventry{Embedded Linux Systems Intern --- ODC}{Dec 2023 -- Jan 2024}
\begin{itemize}
  \item Built Linux systems with Yocto and Buildroot; worked with kernel modules, device drivers, device trees, file systems and hardware flashing.
\end{itemize}

\section*{Technical Skills}
\cvskill{Languages}{Rust, C, C++ (modern C++, Qt), Python, Bash, Java}
\cvskill{Systems \& Linux}{Linux (Ubuntu/Debian), operating systems, FFI and bindings, IPC (varlink, D-Bus; zbus, zlink), concurrency (async/event-driven, multithreading), memory management, networking (TCP/UDP), CUPS/IPP, DNS-SD/mDNS (Avahi), kernel modules and device drivers}
\cvskill{Tooling}{Cargo, Clippy, rustfmt, Git/GitHub, CI, CMake, Make, Yocto, Buildroot, QEMU, testing, debugging, performance optimization, code review, documentation}
\cvskill{Desktop \& embedded}{iced, libcosmic, Qt; Zephyr, FreeRTOS, ARM Cortex-M, ESP32}

\section*{Achievements}
\begin{itemize}
  \item Codeforces Expert (handle \_Khalifa) and ECPC 2023 finalist (Egyptian Collegiate Programming Contest).
  \item Deloitte Software Engineering Mentorship (React Native, 2024); IEEE student chapter Treasurer (\$10k+ budget).
\end{itemize}

\section*{Education}
\cventry{B.Sc. Computer and Systems Engineering --- Minia University}{2020 -- 2025}
Very Good (84.6\%) with Honors. Graduation project: AUTOSAR-based V2X vehicle prototype.

\end{document}
